\documentclass[11pt,a4paper]{article}
\usepackage[margin=26mm]{geometry}
\usepackage[T1]{fontenc}
\usepackage{lmodern,amsmath,amssymb,amsthm,microtype,xcolor}
\setlength{\emergencystretch}{1em}
\usepackage[colorlinks=true,linkcolor=blue!45!black,urlcolor=blue!45!black,citecolor=blue!45!black]{hyperref}
\newtheorem{theorem}{Theorem}
\newtheorem{lemma}{Lemma}
\newcommand{\Q}{\mathbb Q}
\title{Algebraic spherical configurations\\and the Euclidean Ramsey tensor criterion}
\author{AI-assisted research draft\thanks{AI assistance: OpenAI ChatGPT was used for mathematical exploration and drafting; OpenAI Codex was used for editing and compilation.}}
\date{8 October 2026}
\hypersetup{pdfsubject={Research draft}}
\begin{document}
\maketitle
\begin{abstract}
Assuming the tensor classification in OpenAI/math family 172, every finite nonempty spherical configuration with real algebraic coordinates is Euclidean Ramsey. The deduction uses the separability idempotent of its coordinate field. This note records the short argument in full. It does not independently verify the source classification or assert priority for the corollary.
\end{abstract}

\section{The source criterion}
A nonempty finite Euclidean configuration is Ramsey if, for each finite number of colors, some finite-dimensional Euclidean space has a monochromatic congruent copy of the configuration in every coloring. The scale is fixed, and the colorings need not be measurable.

We use the following statement as an external input: \cite[Theorem 1.1]{source}. Let $A=\{a_1,\ldots,a_s\}\subset\mathbb R^d$ have full affine span, with $s\ge2$. Put
\[
 p_i=(1,a_i)^T,\qquad F=\Q((a_i)_\alpha),\qquad B=F\otimes_\Q F,
 \qquad m(x\otimes y)=xy.
\]
The configuration is Ramsey if and only if there is $P\in\operatorname{Mat}_{d+1}(B)$ with
\begin{align}
 (p_i\otimes1)^TP(1\otimes p_i)&=0\quad(1\le i\le s),\label{evaluation}\\
 m(P_{\alpha\beta})&=\delta_{\alpha\beta}\quad(1\le\alpha,\beta\le d).\label{metric}
\end{align}
The coordinate index $0$ denotes the constant coordinate.

\section{A separability calculation}
\begin{lemma}
For a number field $F$, there exists $e\in F\otimes_\Q F$ such that
\[
 m(e)=1,\qquad e(x\otimes1)=e(1\otimes x)\quad(x\in F).
\]
\end{lemma}
\begin{proof}
Choose a primitive element $\theta$ and its separable minimal polynomial $f$. Under $F\otimes_\Q F\simeq F[T]/(f(T))$, the second copy of $\theta$ corresponds to $T$ and $m$ is evaluation at $T=\theta$. Write $f(T)=(T-\theta)g(T)$. Separability gives $g(\theta)\ne0$. The Chinese remainder decomposition
\[
 F[T]/(f)\simeq F[T]/(T-\theta)\times F[T]/(g)
\]
has an idempotent $e$ corresponding to $(1,0)$. Thus $m(e)=1$ and $e(T-\theta)=0$. Every element of $F$ is a rational polynomial in $\theta$, so the asserted identity follows.
\end{proof}

\begin{theorem}
Assume the source tensor criterion. Every finite nonempty spherical set whose coordinates are real algebraic numbers is Euclidean Ramsey.
\end{theorem}
\begin{proof}
A singleton is immediate. Reduce to the affine span of the configuration. This reduction can be made using an algebraic orthonormal basis: apply Gram--Schmidt to a maximal independent list of differences of points. All divisions and square roots involved remain algebraic. The resulting full-span coordinates therefore generate a number field $F$.

Sphericity supplies a real solution $(\ell,c)$ to
\[
 \|a_i\|^2+\ell\cdot a_i+c=0\quad(1\le i\le s).
\]
This is a consistent linear system over $F$, so Gaussian elimination supplies a solution in $F^{d+1}$. Let
\[
 S=\begin{pmatrix}c&\ell^T/2\\ \ell/2&I_d\end{pmatrix}.
\]
Then $p_i^TSp_i=0$ for every $i$. Take the idempotent $e$ from the lemma, and set $P=e(S\otimes1)$ entrywise. Moving the second tensor coordinates through $e$ gives
\[
 (p_i\otimes1)^TP(1\otimes p_i)
 =e\bigl((p_i^TSp_i)\otimes1\bigr)=0.
\]
Moreover $m(P)=S$, so its spatial block is $I_d$. Equations \eqref{evaluation}--\eqref{metric} now imply the Ramsey property.
\end{proof}

\section{A seven-point algebraic subfamily}
For $t>2$ set
\[
 P_t=\{(0,0)\}\cup
 \{(j,\pm\sqrt{j(2t-j)}):j\in\{1,3,4\}\}.
\]
These seven distinct points lie on the circle $(x-t)^2+y^2=t^2$. If $t$ is algebraic, all coordinates are algebraic, and the theorem proves that $P_t$ is Ramsey.

\begin{thebibliography}{9}
\bibitem{source} OpenAI, \emph{A classification of finite Euclidean Ramsey configurations}, 23 September 2026, family 172, Theorem 1.1. \href{https://github.com/openai/math/blob/adc7f1241b42e322a6451854ab7e4b4c146bf78a/preprints/A-classification-of-finite-Euclidean-Ramsey-configurations-September-23-2026/paper.pdf}{Repository snapshot}.
\end{thebibliography}
\end{document}
